\newcommand{\introduction}{
\softwareDefinition
\openSourceSoftwareDefinition
\openSourceToolsDefinition
}

\newcommand{\softwareDefinition}{\par\indent \textbf{Software} é o conjunto de programas, instruções e dados que permitem ao computador executar diferentes tarefas,\cite{david:2001} \textbf{Software Livre e Ferramentas Open Source} são softwares que tem direitos de cópia e distribuição que garantem a sua livre distribuição, uso, modificação e estudo. São representados pela parte lógica do computador, diferentemente do hardware, que corresponde às partes materiais como, processador, placa de video, placa-mãe e circuitos. O software funciona como um conjunto de comandos que orienta o hardware sobre o que deve ser feito, permitindo que o usuário interaja com a máquina. O software é fundamental para o funcionamento dos dispositivos computacionais, pois possibilita que o hardware execute instruções e que os usuários utilizem as máquinas para diversas finalidades.\par \openSourceSoftwareDefinition}

\newcommand{\openSourceSoftwareDefinition}{Diferentemente dos softwares proprietários\cite{fsf:2023}, em que o código-fonte é fechado e controlado por uma empresa ou desenvolvedor, o software livre permite que qualquer pessoa tenha acesso ao código e possa analisá-lo, adaptá-lo ou melhorá-lo conforme suas necessidades. Esse tipo de modelo incentiva a colaboração entre programadores e comunidades ao redor do mundo, contribuindo para o desenvolvimento e aprimoramento contínuo dos programas. \par Além disso, o software livre está presente em diferentes categorias amplamente utilizadas na computação moderna. Entre os principais tipos, destacam-se os sistemas operacionais, como o Linux, que servem como base para a execução de outros programas, os sistemas de gerenciamento de banco de dados, como o MySQL, as ferramentas de desenvolvimento e controle de versão, como o Git,além de aplicações de uso cotidiano como o navegador Mozilla Firefox. \cite{weber:2004}A  ampla adoção dessas ferramentas demonstra como o modelo \textbf{Open Source} se consolidou como um dos pilares da infraestrutura digital global, permitindo inovação contínua e colaboração em larga escala.}

\newcommand{\openSourceToolsDefinition}{ \par \textbf Esse O modelo \textbf{Open Source} de desenvolvimento se baseia em princípios estabelecidos por organizações como a Open Source Initiative e a Free Software Foundation, que definem as liberdades e as estruturas de licenciamento que sustentam os ecossistemas abertos. O desenvolvimento de código aberto se beneficia de contribuições descentralizadas \cite{raymond:2001, fogel:2005}, permitindo que o software evolua de forma mais eficiente em comparação com modelos fechados, destacando-se a importância da colaboração impulsionada pela comunidade para a sustentabilidade desses projetos. Economicamente, o código aberto reduz custos e incentiva a inovação por meio do compartilhamento de conhecimento \cite{lerner:2002}. Exemplos práticos desse paradigma incluem ferramentas amplamente adotadas, como o sistema operacional Linux e Android, sistemas de bancos de dados relacionais como o PostgreSLQ e MySQL que receberam melhorias de performance por colaboradores e pesquisadores externos \cite{repas:2022},demonstram como a colaboração aberta pode resultar em soluções tecnológicas robustas, escaláveis e confiáveis, utilizadas tanto na academia quanto na indústria.}